\documentclass[11pt,a4paper]{article}

\usepackage[margin=2.4cm]{geometry}
\usepackage{amsmath,amssymb,bm}
\usepackage[T1]{fontenc}
\usepackage[utf8]{inputenc}
\usepackage{natbib}
\usepackage{booktabs}
\usepackage{hyperref}
\hypersetup{colorlinks=true,linkcolor=blue!50!black,citecolor=blue!50!black,urlcolor=blue!50!black}
\usepackage{xcolor}

\bibliographystyle{plainnat}
\setcitestyle{authoryear,round}

\newcommand{\Ndot}{\dot{N}_{\rm ion}}
\newcommand{\aB}{\alpha_{\rm B}}
\newcommand{\aA}{\alpha_{\rm A}}
\newcommand{\nH}{n_{\rm H}}
\newcommand{\HII}{\ion{H}{ii}}
\newcommand{\ion}[2]{\textrm{#1\,\textsc{#2}}}
\newcommand{\cms}{\,{\rm cm^{3}\,s^{-1}}}
\newcommand{\pMpc}{\,{\rm pMpc}}
\newcommand{\cMpc}{\,{\rm cMpc}}

\title{\bf The Str\"omgren Sphere: Formation, Evolution and its Role\\ in the Very Early Universe\\[2mm]
\large A concept-and-equation summary built from the papers in \texttt{references/}}
\author{Compiled for L.\ Carvalho}
\date{2 September 2026}

\begin{document}
\maketitle

\begin{abstract}
\noindent
This note collects the physical concepts and the equations that govern the formation and the
evolution of a Str\"omgren sphere, from the classical stellar \ion{H}{ii} region
\citep{Stromgren1939} to the cosmological \ion{H}{ii} regions (``proximity zones'', ``ionized
bubbles'') that surround the first stars, the first quasars and the JWST-era broad-line sources at
$z\gtrsim6$. Every physical statement below is tied to one of the references stored in
\texttt{references/} and listed in \texttt{references.bib}. Section~\ref{sec:checks} reports the
internal-consistency checks (analytic, with \textsc{sympy}, and numerical) that were run on the
equations quoted here; the scripts are \texttt{check\_consistency.py} and
\texttt{alpha\_HI\_recombination\_fits.py}.
\end{abstract}

\tableofcontents

%=====================================================================
\section{The classical Str\"omgren sphere}
\label{sec:classical}

\subsection{The idea}
\citet{Stromgren1939} showed that the Balmer emission observed by Struve \& Elvey in the diffuse
interstellar medium can be understood if the ionizing radiation of an O-type star is absorbed inside
a well-defined volume: \emph{``Balmer-line emission should be limited to certain rather sharply
bounded regions in space surrounding O-type stars or clusters of O-type stars''}, with diameters of
order $200$~pc for the interstellar densities he assumed \citep{Stromgren1939}. The sharpness of the
boundary is a consequence of the very large photoionization cross-section at the Lyman limit: the
transition layer is optically thin in size compared to the ionized volume.

\subsection{Ionization balance and the Str\"omgren radius}
In a static, uniform, pure-hydrogen medium, a steady source emitting $\Ndot$ hydrogen-ionizing
photons per second builds an ionized sphere whose radius follows from the requirement that
\emph{every ionizing photon is balanced by one recombination} inside the sphere
\citep[eq.~1]{Cen2000}:
\begin{equation}
  R_{\rm S} \;=\; \left[\frac{3\,\Ndot}{4\pi\,\aB\,\langle n_{\rm H}^2\rangle}\right]^{1/3},
  \label{eq:RS}
\end{equation}
where $\aB$ is the hydrogen recombination coefficient at $T=10^4$~K and $\langle n_{\rm H}^2\rangle$
is the mean squared hydrogen density inside $R_{\rm S}$ \citep{Cen2000}. Equation~(\ref{eq:RS}) is
the defining relation of the Str\"omgren sphere.

\subsection{Case A, Case B and the on-the-spot approximation}
The choice of recombination coefficient is not cosmetic. Recombinations directly to the ground state
($n=1$) produce a fresh Lyman-continuum photon. The \emph{on-the-spot} (OTS) approximation assumes
that such a photon is re-absorbed close to where it was emitted, so that ground-state recombinations
can simply be omitted; one then uses Case~B, $\aB$, the sum over recombinations to all levels
\emph{except} the ground state \citep{Nebrin2023}. If instead the diffuse photons escape, the correct
coefficient is Case~A, $\aA$, the sum over all levels \citep{Nebrin2023}. Standard values are
\begin{align}
  \aB(10^4\,{\rm K}) &= 2.59\times10^{-13}\cms \quad \text{\citep{Shapiro2006}},\\
  \aA(T) &\simeq 4.2\times10^{-13}\,T_4^{-0.7}\cms \quad \text{\citep{Gnedin2022,Choudhury2022}},
\end{align}
with $T_4\equiv T/10^4$~K. \citet{Nebrin2023} derived an effective coefficient that interpolates
between the two limits for a cloud of Lyman-limit optical thickness $\tau_{\rm cl}$,
\begin{equation}
  \langle\alpha_{\rm eff}\rangle \;=\; \aA - (\aA-\aB)\,f(\tau_{\rm cl}),
  \label{eq:aeff}
\end{equation}
with $f(\tau_{\rm cl})\ll1$ (hence $\langle\alpha_{\rm eff}\rangle\simeq\aA$) for optically thin
clouds and $f(\tau_{\rm cl})\simeq1$ (hence $\langle\alpha_{\rm eff}\rangle\simeq\aB$) for optically
thick clouds \citep{Nebrin2023}. Adopting Case~A where Case~B applies (or vice versa) without
solving for the diffuse field introduces systematic errors in the predicted ionization state
\citep{Nebrin2023}.

\subsection{Temperature dependence of the H\,\textsc{i} recombination coefficient}
\label{sec:alphaT}
Equation~(\ref{eq:RS}) is usually evaluated at the fiducial $T=10^4$~K, but $\aB$ is a steep
function of temperature, and Sec.~\ref{sec:thermal} shows that the gas behind a cosmological I-front
can be anywhere between $\sim1.7\times10^4$ and $3\times10^4$~K. The relevant expressions are
collected here; all of them are fits to detailed collisional-radiative calculations, not independent
theory. They are implemented in \texttt{alpha\_HI\_recombination\_fits.py}.

\paragraph{Case B (recommended closed form).}
\citet{Hui1997} give, in their Appendix~A, a fit to the data of Ferland et al.\ (1992) that is
quoted as accurate to $0.7\%$ over $1\,{\rm K}\le T\le10^9$~K:
\begin{equation}
  \boxed{\;
  \aB(T) \;=\; 2.753\times10^{-14}\,
  \frac{\lambda^{1.500}}{\bigl[1+(\lambda/2.740)^{0.407}\bigr]^{2.242}}\ \cms\;}
  \qquad
  \lambda \;\equiv\; \frac{2\,T_{\rm H\,\textsc{i}}}{T},
  \label{eq:aB_hui}
\end{equation}
with $T_{\rm H\,\textsc{i}}=157807$~K the H\,\textsc{i} ionization threshold expressed in
temperature units \citep{Hui1997}. This is the single most convenient expression for
Str\"omgren-sphere work: it is analytic, has the correct asymptotic behaviour, and is smooth across
the whole nebular and IGM range.

\paragraph{Case A, same source.}
\begin{equation}
  \aA(T) \;=\; 1.269\times10^{-13}\,
  \frac{\lambda^{1.503}}{\bigl[1+(\lambda/0.522)^{0.470}\bigr]^{1.923}}\ \cms,
  \label{eq:aA_hui}
\end{equation}
quoted as accurate to $2\%$ over $3\,{\rm K}\le T\le10^9$~K \citep{Hui1997}.

\paragraph{Total (Case A) rate from atomic-data fits.}
\citet{Verner1996} fitted the total radiative recombination rate of all H-like ions over
$3\,{\rm K}$--$10^{10}$~K with
\begin{equation}
  \alpha_r(T) \;=\; a\left[\sqrt{T/T_0}\,
  \Bigl(1+\sqrt{T/T_0}\Bigr)^{1-b}
  \Bigl(1+\sqrt{T/T_1}\Bigr)^{1+b}\right]^{-1},
  \label{eq:verner}
\end{equation}
which by construction gives $\alpha_r\propto T^{-1/2}$ for $T\ll T_0\ll T_1$ and
$\alpha_r\propto T^{-3/2}$ for $T\gg T_1\gg T_0$ \citep{Verner1996}. Their Table~1 lists, for
H\,\textsc{i},
\begin{equation}
  a = 7.982\times10^{-11}\cms,\quad b=0.7480,\quad T_0=3.148\ {\rm K},\quad T_1=7.036\times10^{5}\ {\rm K},
  \label{eq:verner_HI}
\end{equation}
with rms error $<0.5\%$ (maximum $<1\%$) for hydrogenic species \citep{Verner1996}. Hydrogenic rates
for any charge follow from $\alpha_r(Z,T)=Z\,\alpha_r(1,T/Z^2)$ \citep{Verner1996}. The same
functional form, with $b\to b+b_1\exp(T_2/T)$ for some ions, is what \citet{Mao2016} adopt for the
level-resolved rates, so eq.~(\ref{eq:verner}) is the current standard shape.

\paragraph{Older analytic forms, for provenance.}
Two expressions still widely encountered are, first, the hydrogenic formula of Seaton (1959) as
quoted in eq.~(6) of \citet{Verner1996},
\begin{equation}
  \alpha_r(Z,T) \;=\; 5.197\times10^{-14}\,Z\,\lambda_{\rm S}^{1/2}
  \left[0.4288 + \tfrac{1}{2}\ln\lambda_{\rm S} + 0.469\,\lambda_{\rm S}^{-1/3}\right]\cms,
  \qquad \lambda_{\rm S}=\frac{157890\,Z^2}{T({\rm K})},
  \label{eq:seaton}
\end{equation}
which \citet{Verner1996} note is \emph{not} valid for $T>10^6Z^2$~K because its high-temperature
asymptote is wrong; and second, the four-parameter form of P\'equignot, Petitjean \& Boisson (1991),
also quoted by \citet{Verner1996},
\begin{equation}
  \alpha_r(T) \;=\; 10^{-13}\,z\,\frac{a\,t^{\,b}}{1+c\,t^{\,d}}\cms,
  \qquad t = 10^{-4}T/z^2,
  \label{eq:pequignot}
\end{equation}
valid only over $40z^2\le T\le4\times10^4z^2$~K \citep{Verner1996}.

\paragraph{Polynomial form used in primordial-chemistry networks.}
\citet{Abel1997} fit the same Ferland et al.\ (1992) data with a ninth-order polynomial in
$\ln T$,
\begin{equation}
  k_2 \;=\; \exp\!\left[\sum_{i=0}^{9} c_i \,(\ln T)^i\right]\cms,
  \label{eq:abel}
\end{equation}
with $c_0=-28.6130338$, $c_1=-0.72411256$, $c_2=-2.02604473\times10^{-2}$,
$c_3=-2.38086188\times10^{-3}$, $c_4=-3.21260521\times10^{-4}$, $c_5=-1.42150291\times10^{-5}$,
$c_6=+4.98910892\times10^{-6}$, $c_7=+5.75561414\times10^{-7}$, $c_8=-1.85676704\times10^{-8}$,
$c_9=-3.07113524\times10^{-9}$ \citep{Abel1997}.
\textbf{Unit trap:} the caption of their Table~1 states that ``the temperatures are in eV unless
stated otherwise'', so the argument of eq.~(\ref{eq:abel}) is $\ln(T/{\rm eV})$, not $\ln(T/{\rm
K})$. Using kelvin gives a value $\sim4$ orders of magnitude too small; with $T$ in eV,
eq.~(\ref{eq:abel}) reproduces eq.~(\ref{eq:aA_hui}) to within $3\%$ over $10^2$--$10^6$~K, i.e.\ it
is a \emph{Case A} rate (see check~10 in Sec.~\ref{sec:checks}).

\paragraph{Power-law approximations and their validity.}
Differentiating eq.~(\ref{eq:aB_hui}) numerically gives a local logarithmic slope
${\rm d}\ln\aB/{\rm d}\ln T = -0.787$, $-0.834$ and $-0.888$ at $T=5\times10^3$, $10^4$ and
$2\times10^4$~K, and for Case A ${\rm d}\ln\aA/{\rm d}\ln T = -0.685$, $-0.714$ and $-0.751$ at the
same temperatures. The power laws
\begin{align}
  \aB(T) &\simeq 2.59\times10^{-13}\,T_4^{-0.83}\cms, \label{eq:plB}\\
  \aA(T) &\simeq 4.2\times10^{-13}\,T_4^{-0.70}\cms, \label{eq:plA}
\end{align}
are accurate to $2.1\%$ and $2.5\%$ respectively over $5\times10^3$--$2\times10^4$~K, degrading to
$5.2\%$ and $2.5\%$ over $5\times10^3$--$3\times10^4$~K and to $25\%$ and $19\%$ if stretched over
$10^3$--$10^5$~K. Equation~(\ref{eq:plA}) is exactly the expression quoted by \citet{Gnedin2022} and
\citet{Choudhury2022}; there is \emph{no} single power law that is simultaneously accurate for
Case~A and Case~B, because the two have different slopes. Since $R_{\rm S}\propto\aB^{-1/3}$
(eq.~\ref{eq:RS}), an error $\delta$ in $\aB$ propagates to only $\delta/3$ in $R_{\rm S}$, so
eqs.~(\ref{eq:plB})--(\ref{eq:plA}) are adequate for radii; they are not adequate for photon-budget
or emission-measure work, where $\aB$ enters linearly.

\paragraph{Regime of validity, and what these fits do \emph{not} include.}
All the expressions above are for pure radiative recombination in the low-density limit, with a
Maxwell--Boltzmann electron energy distribution. Two caveats matter for early-Universe work.
(i)~Case~B is itself density-dependent through the collisional redistribution of the excited levels:
\citet{Storey2014} tabulate the hydrogen total Case~B coefficient as $\alpha(N_e,T_e,\kappa)$, i.e.\
explicitly as a function of electron density as well as temperature, so eq.~(\ref{eq:aB_hui}) should
be read as the low-density limit.
(ii)~If the electron energy distribution departs from Maxwellian --- as it may where non-thermal
particles inject energy --- the recombination coefficients change: \citet{Storey2014} compute them
for a $\kappa$-distribution, recovering the Maxwell--Boltzmann values only as $\kappa\to\infty$.
Neither effect is captured by eqs.~(\ref{eq:aB_hui})--(\ref{eq:plA}).

\subsection{Time-dependent growth: the recombination clock}
Before equilibrium is reached, the ionization front (I-front) sweeps up neutral gas. Writing the
photon budget as a balance between injection and recombination inside the ionized volume gives, for a
static uniform medium of clumping factor $C$,
\begin{equation}
  \frac{{\rm d}R^3}{{\rm d}t} \;=\; \frac{3\,\Ndot}{4\pi\,\nH} \;-\; \aB\,C\,\nH\,R^3,
  \label{eq:ode_static}
\end{equation}
whose solution is the familiar exponential approach to the Str\"omgren radius,
\begin{equation}
  R(t) \;=\; R_{\rm S}\,\bigl(1-e^{-t/t_{\rm rec}}\bigr)^{1/3},
  \qquad
  t_{\rm rec} \;=\; \bigl(C\,\aB\,\nH\,\chi_{\rm eff}\bigr)^{-1},
  \label{eq:Rt}
\end{equation}
in agreement with the non-dimensional variable $w\equiv t/t_{\rm rec}=t\,C\,\aB\,\nH\,\chi_{\rm eff}$
and the solution $y=1-e^{-w}$ of \citet{Shapiro2006}; $\chi_{\rm eff}$ accounts for the extra free
electrons contributed by ionized helium. The recombination time is therefore the clock of the
problem: the sphere is essentially complete after a few $t_{\rm rec}$.

\subsection{Relativistic (finite light-speed) corrections}
\citet{Shapiro2006} showed that the non-relativistic equation~(\ref{eq:Rt}) can be badly wrong at
early times, because the I-front is initially so fast that a naive treatment yields apparent
superluminal speeds. The controlling dimensionless parameter is
\begin{equation}
  q \;\equiv\; \frac{R_{\rm S}}{c\,t_{\rm rec}},
  \label{eq:q}
\end{equation}
the light-crossing time of the Str\"omgren radius in units of the recombination time
\citep{Shapiro2006}. For $q>1$ the expansion is relativistic all the way out to $R_{\rm S}$ and the
time to reach the Str\"omgren radius is delayed by a factor $\sim q$ relative to the non-relativistic
answer \citep{Shapiro2006}. This is the regime relevant to bright cosmological sources, where
$t_{\rm rec}$ is enormous (Sec.~\ref{sec:cosmo}).

%=====================================================================
\section{Ionization fronts: R-type, D-type, and the dynamical phase}
\label{sec:fronts}

\subsection{Front classification}
Around a newly switched-on source the I-front first moves supersonically with respect to both the
neutral and the ionized gas: it is a weak \emph{R-type} front, and the gas is ionized essentially
without being moved \citep{Shapiro2006}. As the front decelerates (its speed drops roughly as the
volume grows) it approaches the sound speed of the ionized gas, $c_i$, and transforms into a subsonic
\emph{D-type} front preceded by a shock; from then on the hot, over-pressured ionized bubble expands
hydrodynamically into the surrounding medium \citep{Shapiro2006}. The R-type phase is short-lived
(at most a few $t_{\rm rec}$), whereas the D-type phase is long-lived, and during it $v_{\rm
I}\ll c$ so that relativistic corrections are negligible \citep{Shapiro2006}.

\subsection{D-type expansion: Spitzer and Hosokawa--Inutsuka solutions}
In the thin-shell approximation, equating the pressure of the shocked neutral shell to the ram
pressure of the undisturbed gas gives the ``Raga-I'' equation, whose early-time limit is the
classical \emph{Spitzer solution} \citep[eq.~9]{Bisbas2015}:
\begin{equation}
  R_{\rm Sp}(t) \;=\; R_{\rm St}\left(1+\frac{7}{4}\,\frac{c_i\,t}{R_{\rm St}}\right)^{4/7},
  \label{eq:spitzer}
\end{equation}
with $R_{\rm St}$ the initial Str\"omgren radius of eq.~(\ref{eq:RS}). Including the inertia of the
shocked shell (Hosokawa \& Inutsuka; ``Raga-II'') yields
\begin{equation}
  \dot R_{\rm HI}(t) \;=\; c_i\sqrt{\frac{4}{3}\left(\frac{R_{\rm St}}{R_{\rm HI}}\right)^{3/2}
  -\frac{\mu_i T_o}{\mu_o T_i}},
  \label{eq:HI}
\end{equation}
\citep[eq.~12]{Bisbas2015}. Expansion stops when the ionized gas reaches pressure equilibrium with
the ambient medium, at the \emph{stagnation radius}
\begin{equation}
  R_{\rm STAG} \;=\; \left(\frac{c_i}{c_o}\right)^{4/3} R_{\rm St},
  \label{eq:stag}
\end{equation}
where $c_o$ is the sound speed of the neutral gas \citep[eqs.~14 and 22]{Bisbas2015}. Because
$c_i/c_o\sim4$--$10$, the final \ion{H}{ii} region is one to two orders of magnitude larger in volume
than the initial Str\"omgren sphere. \citet{Bisbas2015} showed with the \textsc{StarBench} code
comparison that a semi-empirical combination of the early- and late-time solutions reproduces
high-resolution radiation-hydrodynamics simulations to $\lesssim2\%$, and is a better code-validation
benchmark than the Spitzer solution alone.

\subsection{The first \ion{H}{ii} regions: Population~III stars in minihaloes}
In the primordial case the ambient density profile is steeply declining, and the ordering of the two
front types is inverted. \citet{Kitayama2004} found that the \ion{H}{ii} region of a Pop~III star
begins as a slowly expanding weak D-type front near the centre and then, once the front outruns the
declining density profile, converts to a rapidly propagating R-type front through the outer envelope;
this D$\to$R transition is the critical condition for complete ionization of a cosmologically
relevant halo. For halo masses $\lesssim10^6\,M_\odot$ the transition occurs within a few
$10^5$~yr, the escape fractions of both ionizing and photodissociating photons exceed $80\%$, and a
supersonic shock evacuates the halo, lowering the mean density to $\lesssim1\,{\rm cm^{-3}}$ in a few
Myr; for $\gtrsim10^7\,M_\odot$ haloes the front stays D-type for the whole stellar lifetime, the
\ion{H}{ii} region remains confined inside the virial radius and the escape fraction is essentially
zero \citep{Kitayama2004}. \citet{Alvarez2006}, with 3D ray tracing, described the same physics as a
``champagne flow'' in which the early D-type front detaches from the shock and runs ahead as an
R-type front; they obtained $0.7\lesssim f_{\rm esc}\lesssim0.9$ for $80\lesssim M_\star/M_\odot
\lesssim500$ and a ratio of ionized gas mass to stellar mass of $\sim6\times10^4$, i.e.\ roughly half
the ionizing photons released per stellar baryon are used to ionize the intergalactic gas. They also
found that neighbouring minihaloes \emph{trap} the I-front, so that their cores remain neutral
\citep{Alvarez2006}.

%=====================================================================
\section{Cosmological Str\"omgren spheres}
\label{sec:cosmo}

\subsection{Equation of motion in an expanding universe}
For a cosmological source the Str\"omgren problem acquires two new ingredients: the Hubble expansion
dilutes the gas, and the intergalactic medium (IGM) is clumpy. Following
\citet{ShapiroGiroux1987}, the proper radius $R_i$ of the I-front obeys \citep[eq.~3]{Cen2000}
\begin{equation}
  \frac{{\rm d}R_i^3}{{\rm d}t}
  \;=\; 3H(z)\,R_i^3
  \;+\; \frac{3\,\dot N_{\rm ph}}{4\pi\,\langle \nH\rangle}
  \;-\; C_{\rm HII}\,\langle \nH\rangle\,\aB\,R_i^3,
  \label{eq:cosmo_ode}
\end{equation}
where the three terms on the right-hand side are, in order, the Hubble expansion, the ionizations by
newly produced photons, and the recombinations, and
\begin{equation}
  C_{\rm HII}\equiv\frac{\langle n_{\rm H}^2\rangle}{\langle \nH\rangle^2}
\end{equation}
is the mean clumping factor of the ionized gas inside $R_i$ \citep{Cen2000}. A more careful
definition weights by the electron density,
$C\equiv\langle n_{\rm HII}n_e\rangle/(\bar n_{\rm HII}\bar n_e)$, restricted to the ionized regions
\citep[eq.~64]{Choudhury2022}; $C=1$ for a uniform medium. The clumping factor is the standard way of
encoding IGM inhomogeneity as an enhanced effective recombination rate
\citep{MiraldaEscude2000,Choudhury2022}, and it is not straightforward to model self-consistently
because it depends on the density structure and on how I-fronts propagate into dense regions
\citep{Choudhury2022}.

\subsection{The photon-counting limit: why high-$z$ spheres never reach equilibrium}
The key qualitative difference with respect to the stellar case is that at $z\lesssim10$ the
\emph{recombination time of the IGM is comparable to or longer than the Hubble time, and enormously
longer than the quasar lifetime}, so the equilibrium radius $R_{\rm S}$ is large but is never
actually reached before the source turns off \citep{Cen2000}. (For $C=1$ we find
$t_{\rm rec}/H^{-1}(z)\simeq1.2$, $0.8$ and $0.6$ at $z=6$, $8$ and $10$; see check~7 in
Sec.~\ref{sec:checks}. The statement is therefore safe with respect to $t_Q$, which is what matters
here, and marginal with respect to the Hubble time at the highest redshifts.) Dropping the recombination term in eq.~(\ref{eq:cosmo_ode}) leaves a pure
photon-counting result for a source of age $t_Q$ \citep[eq.~2]{Cen2000}:
\begin{equation}
  R_{t_Q} \;=\; \left[\frac{3\,\dot N_{\rm ph}\,t_Q}{4\pi\,\langle \nH\rangle}\right]^{1/3}.
  \label{eq:photon_counting}
\end{equation}
Including the neutral fraction $x_{\rm HI}$ of the gas that must be ionized, the fiducial scaling for
a luminous $z\sim6$ quasar is \citep{Mesinger2004}
\begin{equation}
  R_{\rm S}\simeq 7.7\;
  x_{\rm HI}^{-1/3}
  \left(\frac{\dot N_Q}{6.5\times10^{57}\,{\rm s^{-1}}}\right)^{1/3}
  \left(\frac{t_Q}{2\times10^{7}\,{\rm yr}}\right)^{1/3}
  \left(\frac{1+z_Q}{7.28}\right)^{-1}\ {\rm Mpc\ (proper)}.
  \label{eq:mesinger}
\end{equation}
The three exponents encode the whole physics: $R\propto \dot N^{1/3}$ and $R\propto t_Q^{1/3}$ from
photon counting, and $R\propto(1+z)^{-1}$ from $\langle \nH\rangle\propto(1+z)^3$. Unlike the stellar
case, the cosmological sphere keeps evolving even at constant luminosity, because of the Hubble
expansion, of the evolving density and clumping around the source, and of the continued accumulation
of ionizing photons \citep{Cen2000}.

\subsection{From individual spheres to reionization}
Globally, the same photon budget written for the volume filling factor $Q_{\rm HII}$ of ionized gas
gives the standard reionization equation \citep[eq.~66]{Choudhury2022}
\begin{equation}
  \frac{{\rm d}Q_{\rm HII}}{{\rm d}t}
  \;=\; \frac{\dot n_\gamma}{\bar n_{\rm H}}
  \;-\; Q_{\rm HII}\,a^{-3}\,C\,\alpha_{\rm R}(T)\,\bar n_e ,
  \label{eq:QHII}
\end{equation}
i.e.\ a source term set by the ionizing emissivity and a sink term set by the clumping-enhanced
recombination rate; $\chi_{\rm He}\approx1.08$ ($1.16$) for singly (doubly) ionized helium with
$Y=0.24$ \citep{Choudhury2022}. Reionization proceeds inside-out: isolated ionized bubbles form
around the first sources, then grow, overlap and finally fill the IGM, leaving neutral gas only in
dense pockets \citep{Chakraborty2025,Cen2000}.

The size \emph{distribution} of these bubbles is not set by the Str\"omgren radius of individual
sources but by the clustering of the sources. \citet{Furlanetto2004} associated each ionized region
with a large-scale density fluctuation that can ``self-ionize'', $m_{\rm ion}=\zeta m_{\rm gal}$,
i.e.
\begin{equation}
  f_{\rm coll}(\delta_m,m,z)\;\ge\;\zeta^{-1},
  \label{eq:barrier}
\end{equation}
and used the excursion-set formalism with the resulting density barrier
$\delta_x(m,z)=\delta_c(z)-\sqrt{2}\,K(\zeta)\,[\sigma^2_{\min}-\sigma^2(m)]^{1/2}$, with
$K(\zeta)={\rm erf}^{-1}(1-\zeta^{-1})$, to obtain the bubble size distribution
\citep{Furlanetto2004}. In this model bubbles are already several comoving Mpc across by the time the
ionized fraction reaches $\bar Q=0.5$ \citep{Furlanetto2004}. Modern radiation-hydrodynamic
simulations confirm that bubble sizes correlate with the local overdensity, and that the local
reionization redshift $z_{\rm reion}$ is a good proxy for bubble sizes only up to
$R_{\rm eff}\lesssim1\cMpc$ \citep{Neyer2023}.

%=====================================================================
\section{Thermal structure: the temperature left behind by the I-front}
\label{sec:thermal}

The Str\"omgren sphere is not only an ionization structure but also a heating structure: the excess
energy of each ionizing photon above $13.6$~eV is deposited as heat when the front passes. The
post-front temperature $T_{\rm reion}$ is a key, and uncertain, ingredient of reionization models
\citep{DAloisio2018}. Its main determinant is \emph{the I-front speed}, not the spectral index of the
source \citep{DAloisio2018,Zhu2025}:
\begin{itemize}
  \item slow fronts ($v_{\rm I}\sim50$--$2000\,{\rm km\,s^{-1}}$) give
        $T\sim17{,}000$--$22{,}000$~K, typical of the first half of reionization
        \citep{DAloisio2018,Zhu2025};
  \item fast fronts near the overlap phase ($v_{\rm I}\sim10^4\,{\rm km\,s^{-1}}$) reach
        $T_{\rm reion}\sim25{,}000$--$30{,}000$~K \citep{DAloisio2018,Zhu2025};
  \item the median post-ionization temperature plateaus at $T\sim2\times10^4$~K for
        $v_{\rm I}>3000\,{\rm km\,s^{-1}}$ \citep{Zhu2025}.
\end{itemize}
The usual assumption that all baryonic species (${\rm e^-}$, \ion{H}{i}, \ion{H}{ii}, \ion{He}{i},
\ion{He}{ii}) instantaneously share the deposited energy breaks down for fast fronts:
\citet{Zeng2020} solved the non-equilibrium energy-transfer problem and found that equilibration is a
good approximation for $v_{\rm I}\lesssim10^9\,{\rm cm\,s^{-1}}$, while $T_{\rm re}$ is lower than
the equilibrium value by up to $19.7\%$ at $3\times10^9\,{\rm cm\,s^{-1}}$ and $30.8\%$ at
$10^{10}\,{\rm cm\,s^{-1}}$. Since electrons carry the photoelectron energy first and then share it
with the ions, the electron--ion energy transfer rate is what sets this deviation \citep{Zeng2020}.
Uncertainties in $T_{\rm reion}$ of $15{,}000$ vs.\ $25{,}000$~K propagate into $\sim50\%$
differences in the IGM temperature at $z\simeq5.5$ for gas reionized at $z_{\rm re}=9$
\citep{DAloisio2018}, which is directly relevant to $z>5$ Ly$\alpha$ forest thermal measurements.

%=====================================================================
\section{Observational manifestations at high redshift}
\label{sec:obs}

\subsection{The Str\"omgren surface in quasar spectra}
A bright quasar embedded in a largely neutral IGM at $z\gtrsim6$ is surrounded by a large
cosmological Str\"omgren sphere; its spectrum then shows a sharp increase in resonant Lyman-line
absorption at wavelengths approaching the boundary of the sphere along the line of sight
\citep{Mesinger2004}. Using \emph{simultaneously} the Ly$\alpha$ and Ly$\beta$ regions to gain
dynamic range, \citet{Mesinger2004} detected the boundary of the Str\"omgren sphere of
SDSS~J1030$+$0524 ($z=6.28$) at a proper distance of $6.0\pm0.2$~Mpc, inferred $x_{\rm HI}\gtrsim0.2$
beyond it, and measured an ionizing output of $(5.2\pm2.5)\times10^{56}$ photons~s$^{-1}$. The
underlying spectra of the two first Gunn--Peterson-trough quasars, and the fact that the Str\"omgren
sphere of SDSS~J1148$+$5251 is significantly smaller, were presented by \citet{White2003}.

The Lyman-$\alpha$ line profile itself carries information: for a source inside a neutral IGM the
damping wing of the intervening \ion{H}{i} suppresses the line, whereas a sufficiently large
\ion{H}{ii} region lets Ly$\alpha$ photons redshift out of resonance before reaching the boundary
\citep{Cen2000}. The red side of the line then measures the quasar lifetime and the blue side the
IGM density fluctuations at the quasar redshift \citep{Cen2000}.

\subsection{Proximity zones and quasar lifetimes}
Operationally, the observed quantity is the \emph{proximity zone} size $R_p$, defined as the distance
at which the smoothed, continuum-normalized quasar spectrum first drops below the $10\%$ transmitted
flux level \citep{Eilers2017}. This definition probes locations where the quasar's ionizing radiation
still exceeds the expected UV background by about an order of magnitude, which makes $R_p$
comparatively \emph{insensitive} to the ambient IGM neutral fraction \citep{Eilers2017}. This is the
main caveat when using near zones as a reionization probe: \citet{Bolton2007} showed with synthetic
spectra that the scatter induced by density fluctuations and by the filtering of ionizing radiation
is so large that the data then available could not distinguish a significantly neutral from a highly
ionized IGM just above $z=6$, and that samples of several tens of high-resolution Ly$\alpha$ +
Ly$\beta$ near zones are needed to decide whether $\langle f_{\rm HI}\rangle_V$ is above or below
$10\%$.

What $R_p$ \emph{does} measure well is the quasar age, because the IGM has a finite response time to
the radiation \citep{Eilers2020}. The observational picture is:
\begin{itemize}
  \item $R_p$ evolves only shallowly with redshift, $\propto(1+z)^{-1.44}$ over
        $5.77\le z\le6.54$, in agreement with radiative-transfer simulations irrespective of the IGM
        ionization state \citep{Eilers2017};
  \item a few quasars have exceptionally small zones ($R_p<1\pMpc$), most compellingly explained by
        UV-luminous lifetimes $t_Q\lesssim10^5$~yr \citep{Eilers2017};
  \item a dedicated multi-wavelength survey found $5\%\lesssim f_{\rm young}\lesssim10\%$ of the
        high-$z$ quasar population to be young, with $t_Q\lesssim10^5$~yr \citep{Eilers2020};
  \item the enlarged XQR-30 sample gives $R_p=2$--$7\pMpc$ for $M_{1450}\sim-27$ quasars at
        $5.8<z<6.6$, with the scatter indicating a large diversity of quasar lifetimes
        \citep{Satyavolu2023};
  \item the largest homogeneous compilation to date, $59$ quasars at $5.77\le z\le7.54$, gives
        $R_p\propto 10^{-0.4M_{1450}/2.87}$ (i.e.\ $R_p\propto L^{1/2.87}$) and
        $R_p\propto(1+z)^{-2.44}$, and identifies $13$ small-zone quasars for which
        $t_Q\lesssim10^4$~yr is the favoured explanation \citep{Onorato2025}.
\end{itemize}
Note that the measured luminosity scaling $L^{1/2.87}=L^{0.348}$ is remarkably close to the naive
photon-counting exponent $1/3$ of eq.~(\ref{eq:photon_counting}), whereas the redshift scaling is
steeper than the $(1+z)^{-1}$ of eq.~(\ref{eq:mesinger}), which is expected since the latter holds at
fixed $x_{\rm HI}$ and fixed $t_Q$ \citep{Onorato2025,Mesinger2004}.

\subsection{Ionized bubbles around JWST galaxies and AGN}
For galaxies, the same physics is inverted: one models the Ly$\alpha$ damping-wing absorption
assuming a galaxy sitting inside an ionized bubble of radius $R_b$ embedded in an IGM of neutral
fraction $x_{\rm HI}$. From $27$ bright ($M_{\rm UV}<-18.5$) JWST/NIRSpec galaxies at $7<z<12$,
\citet{Umeda2023} obtained $x_{\rm HI}$ rising from $0.53^{+0.18}_{-0.47}$ at $z=7.12$ to
$0.92^{+0.08}_{-0.10}$ at $z=9.91$, with $\log(R_b/{\rm cMpc})$ falling from $1.67^{+0.14}_{-0.16}$
to $-0.69^{+1.49}_{-0.24}$ over the same interval; the inferred $R_b$ can be up to a few dex larger
than the cosmic-average analytic expectation at a given $x_{\rm HI}$, but is compatible with the
merged bubbles predicted around bright galaxies in simulations \citep{Umeda2023}. A Ly$\alpha$
line-shape fit that explicitly assumes a Str\"omgren sphere gives $R_b=1.5^{+0.3}_{-0.2}\pMpc$ around
a nitrogen-rich AGN at $z=8.63$ --- a large bubble given that the Universe is nearly fully neutral at
that redshift \citep{Morishita2025}. As quoted by \citet{Morishita2025}, the predicted IGM damping-wing transmission at Ly$\alpha$ line
centre is $<10\%$ if the source sits in an ionized region smaller than $0.5\pMpc$, versus $>30\%$ for
a region larger than $1\pMpc$; this is why the detection of Ly$\alpha$ near the systemic velocity at
$z>8$ is itself evidence for a substantial Str\"omgren sphere.

\subsection{Who blows the bubbles? Relevance for the ionizing photon budget}
Whether the ionizing photons come from faint star-forming galaxies or from accreting black holes
determines both the bubble morphology and the reionization history
\citep{Chakraborty2025,Gnedin2022}. \citet{Mascia2026} carried out a systematic census of compact
broad-line emitters (BLEs) at $4\le z\le7$ from JWST/NIRSpec: the bluest, compact BLEs show hot
ionizing continua, high-ionization UV lines, elevated Ly$\alpha$-emitter fractions
($X_{\rm Ly\alpha}=50$--$67\%$), a high ionizing photon production efficiency
$\log(\xi_{\rm ion}/{\rm Hz\,erg^{-1}})\sim25.4$ and a weighted escape fraction
$f_{\rm esc}\sim0.2$, whereas Little Red Dots have $f_{\rm esc}\sim0.01$. Their conclusion is that
broad-line sources do \emph{not} dominate reionization globally, but that, given their luminosity,
their contribution can dominate \emph{up to few-Mpc scales} \citep{Mascia2026} --- i.e.\ exactly on
the scale of an individual Str\"omgren sphere.

%=====================================================================
\section{Summary of the key equations and fiducial numbers}

\begin{table}[h]
\centering
\small
\begin{tabular}{@{}p{5.6cm}p{6.4cm}p{3.2cm}@{}}
\toprule
Quantity & Expression & Reference \\
\midrule
Str\"omgren radius (static) & $R_{\rm S}=[3\Ndot/(4\pi\aB\langle n_{\rm H}^2\rangle)]^{1/3}$ & \citet{Cen2000} \\
Case B coefficient & $\aB(10^4{\rm K})=2.59\times10^{-13}\cms$ & \citet{Shapiro2006} \\
Case A coefficient & $\aA\simeq4.2\times10^{-13}T_4^{-0.7}\cms$ & \citet{Gnedin2022} \\
Case B, full $T$ dependence & $\aB=2.753\times10^{-14}\lambda^{1.500}[1+(\lambda/2.740)^{0.407}]^{-2.242}$, $\lambda=2T_{\rm H\,\textsc{i}}/T$ & \citet{Hui1997} \\
Case A, full $T$ dependence & $\aA=1.269\times10^{-13}\lambda^{1.503}[1+(\lambda/0.522)^{0.470}]^{-1.923}$ & \citet{Hui1997} \\
Total rate, H-like ions & $\alpha_r=a[\sqrt{T/T_0}(1+\sqrt{T/T_0})^{1-b}(1+\sqrt{T/T_1})^{1+b}]^{-1}$ & \citet{Verner1996} \\
Power-law approximations & $\aB\simeq2.59\times10^{-13}T_4^{-0.83}$, $\aA\simeq4.2\times10^{-13}T_4^{-0.70}$ & Sec.~\ref{sec:alphaT} \\
Effective coefficient & $\langle\alpha_{\rm eff}\rangle=\aA-(\aA-\aB)f(\tau_{\rm cl})$ & \citet{Nebrin2023} \\
Growth to equilibrium & $R(t)=R_{\rm S}(1-e^{-t/t_{\rm rec}})^{1/3}$ & \citet{Shapiro2006} \\
Recombination time & $t_{\rm rec}=(C\aB \nH\chi_{\rm eff})^{-1}$ & \citet{Shapiro2006} \\
Relativistic parameter & $q=R_{\rm S}/(c\,t_{\rm rec})$ & \citet{Shapiro2006} \\
D-type (Spitzer) & $R_{\rm Sp}=R_{\rm St}(1+\tfrac{7}{4}c_it/R_{\rm St})^{4/7}$ & \citet{Bisbas2015} \\
Stagnation radius & $R_{\rm STAG}=(c_i/c_o)^{4/3}R_{\rm St}$ & \citet{Bisbas2015} \\
Cosmological I-front & $\dot{(R_i^3)}=3HR_i^3+\tfrac{3\dot N_{\rm ph}}{4\pi\langle \nH\rangle}-C_{\rm HII}\langle \nH\rangle\aB R_i^3$ & \citet{Cen2000,ShapiroGiroux1987} \\
Photon-counting radius & $R_{t_Q}=[3\dot N_{\rm ph}t_Q/(4\pi\langle \nH\rangle)]^{1/3}$ & \citet{Cen2000} \\
Fiducial high-$z$ scaling & $R_{\rm S}\simeq7.7\,x_{\rm HI}^{-1/3}\dot N_{57}^{1/3}t_{Q,7}^{1/3}(1+z)^{-1}$~pMpc & \citet{Mesinger2004} \\
Reionization equation & ${\rm d}Q_{\rm HII}/{\rm d}t=\dot n_\gamma/\bar n_{\rm H}-Q_{\rm HII}a^{-3}C\alpha_{\rm R}\bar n_e$ & \citet{Choudhury2022} \\
Self-ionization barrier & $f_{\rm coll}\ge\zeta^{-1}$ & \citet{Furlanetto2004} \\
Post-front temperature & $T_{\rm reion}\simeq1.7$--$3.0\times10^4$~K, set by $v_{\rm I}$ & \citet{DAloisio2018,Zhu2025} \\
Proximity zone definition & first drop below $10\%$ transmitted flux & \citet{Eilers2017} \\
\bottomrule
\end{tabular}
\caption{Master list of the relations discussed in this note.}
\end{table}

%=====================================================================
\section{Internal consistency checks}
\label{sec:checks}

The following checks were carried out with \textsc{sympy}~1.12 and \textsc{numpy}~1.26
(\texttt{check\_consistency.py}); all of them passed.

\begin{enumerate}
\item \textbf{Eq.~(\ref{eq:Rt}) solves eq.~(\ref{eq:ode_static}).} Substituting
      $R(t)=R_{\rm S}(1-e^{-\aB C\nH t})^{1/3}$ with
      $R_{\rm S}=[3\Ndot/(4\pi\aB C\nH^2)]^{1/3}$ into
      ${\rm d}R^3/{\rm d}t=3\Ndot/(4\pi \nH)-\aB C\nH R^3$ gives a residual that \textsc{sympy}
      simplifies to exactly $0$. This also confirms the identification
      $t_{\rm rec}=(\aB C \nH)^{-1}$ used in eq.~(\ref{eq:Rt}) and by \citet{Shapiro2006}.
\item \textbf{Eq.~(\ref{eq:spitzer}) solves the early-time Raga-I equation.}
      $\dot R_{\rm Sp}/c_i-(R_{\rm St}/R_{\rm Sp})^{3/4}$ simplifies to $0$, i.e.\ the Spitzer
      solution is the exact solution of eq.~(8) of \citet{Bisbas2015} once the small
      $\mu_iT_o/\mu_oT_i$ term is dropped, as stated there.
\item \textbf{Eq.~(\ref{eq:stag}) reproduces the numbers quoted by \citet{Bisbas2015}.} With
      $T_i=10^4$~K, $\mu_i=0.5$ and $T_o=10^3$~K, $\mu_o=1$ one gets $c_i=12.85$~km\,s$^{-1}$ and
      $c_o=2.87$~km\,s$^{-1}$ (their values), and $R_{\rm STAG}=(c_i/c_o)^{4/3}\times0.314$~pc
      $=2.31$~pc, identical to their quoted stagnation radius.
\item \textbf{Eq.~(\ref{eq:mesinger}) is the photon-counting radius eq.~(\ref{eq:photon_counting}).}
      Evaluating $[3\dot N_Qt_Q/(4\pi \nH x_{\rm HI})]^{1/3}$ with
      $\dot N_Q=6.5\times10^{57}\,{\rm s^{-1}}$, $t_Q=2\times10^7$~yr, $x_{\rm HI}=1$, $z=6.28$ and
      the cosmology of \citet{Mesinger2004} ($\Omega_b=0.044$, $h=0.71$, $Y=0.24$;
      $\nH=7.31\times10^{-5}\,{\rm cm^{-3}}$) gives $7.70$~proper~Mpc, exactly the coefficient
      quoted in their eq.~(\ref{eq:mesinger}).
\item \textbf{Proper/comoving conversion of the measured Str\"omgren surface.} Their best-fit
      $R_{\rm S}=44\cMpc$ at $z=6.28$ corresponds to $44/(1+z)=6.04\pMpc$, consistent with the
      $6.0\pm0.2\pMpc$ quoted in their abstract \citep{Mesinger2004}.
\item \textbf{Order of magnitude of eq.~(\ref{eq:RS}) for a stellar source.} With
      $\Ndot=10^{49}\,{\rm s^{-1}}$ and $\nH=1\,{\rm cm^{-3}}$, eq.~(\ref{eq:RS}) gives
      $R_{\rm S}=68$~pc, i.e.\ a diameter of $136$~pc, of the same order as the ``about 200
      parsecs'' quoted by \citet{Stromgren1939} (the exact value depends on the stellar temperature,
      radius and gas density he adopted). For $\nH=100\,{\rm cm^{-3}}$ the same source gives
      $R_{\rm S}=3.2$~pc and $t_{\rm rec}=1.2$~kyr, so a compact \ion{H}{ii} region reaches
      equilibrium essentially instantaneously --- consistent with the classical, static
      interpretation.
\item \textbf{Why the cosmological sphere is \emph{not} in equilibrium.} With
      $\chi_{\rm He}=1.08$ and $C=1$, $t_{\rm rec}=1.74$, $0.82$ and $0.45$~Gyr at $z=6$, $8$ and
      $10$ respectively, i.e.\ comparable to the Hubble time $1/H(z)=1.43$, $0.98$ and $0.73$~Gyr,
      and four to five orders of magnitude longer than the quasar lifetimes
      $t_Q\sim10^4$--$10^7$~yr inferred from proximity zones
      \citep{Eilers2017,Eilers2020,Onorato2025}. This is a quantitative confirmation of the
      statement of \citet{Cen2000} that the equilibrium Str\"omgren radius is never reached at high
      redshift, and hence that eq.~(\ref{eq:photon_counting}) --- not eq.~(\ref{eq:RS}) --- is the
      relevant relation. Caveat: at $C\gtrsim2$ and $z\gtrsim8$, $t_{\rm rec}$ becomes shorter than
      $1/H(z)$, so the recombination term of eq.~(\ref{eq:cosmo_ode}) cannot be neglected for
      long-lived sources in dense environments.
\item \textbf{Luminosity scaling.} The empirical $R_p\propto L^{1/2.87}=L^{0.348}$ of
      \citet{Onorato2025} is within $5\%$ of the photon-counting exponent $1/3=0.333$ of
      eq.~(\ref{eq:photon_counting}), so the two are mutually consistent.
\item \textbf{I-front speed regimes.} The plateau velocity $3000\,{\rm km\,s^{-1}}
      =3\times10^8\,{\rm cm\,s^{-1}}$ of \citet{Zhu2025} lies \emph{below} the
      $10^9\,{\rm cm\,s^{-1}}$ threshold above which \citet{Zeng2020} find departures from thermal
      equilibration, so the two results are consistent: the $\sim2\times10^4$~K plateau is in the
      regime where the equilibrium assumption holds, and the $19.7$--$30.8\%$ reductions of
      \citet{Zeng2020} apply only to the fastest fronts.
\item \textbf{Mutual consistency of the H\,\textsc{i} recombination fits of
      Sec.~\ref{sec:alphaT}.} Evaluated at $T=10^4$~K, eq.~(\ref{eq:aB_hui}) gives
      $\aB=2.5918\times10^{-13}\cms$, reproducing the $2.59\times10^{-13}\cms$ quoted by
      \citet{Shapiro2006} to $0.07\%$. For Case~A the four independent fits give
      $4.2970\times10^{-13}$ (eq.~\ref{eq:aA_hui}), $4.1923\times10^{-13}$
      (eqs.~\ref{eq:verner}--\ref{eq:verner_HI}), $4.1206\times10^{-13}$ (eq.~\ref{eq:seaton}) and
      $4.1699\times10^{-13}\cms$ (eq.~\ref{eq:abel} with $T$ in eV): they agree to within $4\%$ and
      are all consistent with the $4.2\times10^{-13}\cms$ of \citet{Gnedin2022}. Over
      $10^2$--$10^6$~K the largest deviation between eq.~(\ref{eq:abel}) and eq.~(\ref{eq:aA_hui})
      is $3\%$, which is what identifies eq.~(\ref{eq:abel}) as a Case~A rate. The ratio
      $\aB/\aA=0.603$ at $10^4$~K is the fraction of recombinations landing on excited states, and
      is bracketed by eq.~(\ref{eq:aeff}) as it must be \citep{Nebrin2023}. \emph{Not verified:} the
      underlying Ferland et al.\ (1992) and Storey \& Hummer collisional-radiative calculations
      themselves, which are taken as given by all the fits above.
\end{enumerate}

\paragraph{Known tensions / open issues, stated explicitly.}
(i) The definition of $R_p$ used observationally \citep{Eilers2017} is \emph{not} the Str\"omgren
radius of eq.~(\ref{eq:RS}); it is a flux threshold that traces where the quasar radiation dominates
the background, which is why $R_p$ is a good clock but a poor neutral-fraction meter
\citep{Eilers2017,Bolton2007}.
(ii) The redshift evolution of $R_p$ is measured to be $\propto(1+z)^{-1.44}$ by \citet{Eilers2017}
and $\propto(1+z)^{-2.44}$ by \citet{Onorato2025}; both are steeper than the $(1+z)^{-1}$ expected at
fixed $x_{\rm HI}$ and $t_Q$ (eq.~\ref{eq:mesinger}), and the two measurements differ from each
other, so the redshift scaling should be quoted as observationally uncertain rather than as a
prediction.
(iii) Bubble radii inferred from Ly$\alpha$ damping wings can exceed the analytic cosmic-average
expectation by up to a few dex at a given $x_{\rm HI}$ \citep{Umeda2023}, which is a genuine tension
with single-source Str\"omgren-sphere modelling and points to merged bubbles around clustered sources
\citep{Umeda2023,Neyer2023,Furlanetto2004}.

%=====================================================================
\bibliography{references}

\end{document}
